% Ad Familiares 12.16 --- headnote
% ad-familiares-12-16, 25 May 44 BC, Athens

\textit{Trebonius to Cicero, from Athens on 25 May 44 BC ---
Perseus dateline \textit{Scr. Athenis viii K. Iun. a. 710 (44)}.
This is one of the rare letters in the collection \textbf{from} a
correspondent: C. Trebonius, the consul-suffect of 45, now en route
to govern Asia, writes back to Cicero from Athens just over two
months after the Ides --- in which he had played the (literally
peripheral) part of detaining Antony outside the senate-house door
while the others struck. He reports the one piece of news he knows
Cicero most wants: young Marcus is in Athens, hard at work under the
right teachers, and well-spoken of. The letter is unfussy and warm
in the way only an old friendship licenses.}

\textit{The body has three concerns. First, the report on the son
(with a promise that, when Marcus visits the province, Cratippus
the Peripatetic will be brought along so that the visit is not a
``holiday from his studies''). Second, a literary gift composed at
sea --- a set of verses worked up around one of Cicero's own
\textit{dicta}; Trebonius asks pardon if some of the lines are
\textit{euthyrrh\=emonesteros}, ``a touch more outspoken,'' invoking
Lucilius as licence for satirical freedom. Third, a quiet but
unmissable closing: when Cicero writes \textit{de interitu Caesaris},
he must not leave Trebonius's share of the deed --- or of the
friendship --- out. The famous \textit{quaesivit lucem mortuo
Antonius}: Trebonius would be murdered at Smyrna by Dolabella in
January 43, in what Cicero treats as the first blood of the new war.}
